\makeabstract
{
A Model-independent Radio Telescope Dark Matter Search in the X Band
}
{
Benjamin Hanselman (1,2), Aya Keller (1,3), and Karl van Bibber (1)
}
{
\textsuperscript{1}Department of Nuclear Engineering, University of California, Berkeley\\
\textsuperscript{2}Department of Physics and Astronomy, Amherst College\\
\textsuperscript{3}Department of Physics, Stanford University
}
{
Despite mounting evidence for the existence of dark matter, its identification remains elusive. Recent detector technology advances have enabled experiments that are highly sensitive to specific dark matter candidates like the QCD axion. However, attention has recently been shifting to a theoretical framework for dark matter that does not assume specific models. Our astrophysical dark matter search therefore relies on as few assumptions as possible, focusing on the potential radiative decay or annihilation of ultralight dark matter in the Milky Way halo.
}
{
The X-band search utilized spectroscopic observations of 4999 Hipparcos catalogue targets from the Green Bank Telescope, collected between 2018 and 2025 as part of the Breakthrough Listen campaign. We applied a rolling normalization procedure to each spectrum to remove background variations while preserving potential dark matter signals. We then selected for spectral lines exhibiting the characteristic Doppler and intensity shifts expected from dark matter in the halo by constructing asymmetry spectra. To derive physics limits on the dark matter decay rate λ and velocity-weighted cross-section ⟨σv⟩, we injected synthetic signals into raw spectra, calculated templates for specific dark matter physics parameters, and performed a p-value analysis on the cross-correlation between the asymmetry spectra and the templates.
}
{
We fine-tuned the normalization scheme to maximize the signal-to-noise ratio of intensity asymmetries from synthetic dark matter signals. Optimal results were achieved with a rolling window size of M = 7η and a putative signal blanking region size of m = 1.4η, where η represents the expected dark matter signal width. Our dark matter detection analysis in the 8000-8500 MHz frequency range excluded dark matter decay above λ ≈ 10\textasciicircum{}\{-29\} s\textasciicircum{}\{-1\} and annihilation above ⟨σv⟩ ≈ 10\textasciicircum{}\{-26\} cm\textasciicircum{}3 s\textasciicircum{}\{-1\}.
}
{
Our preliminary exclusion limits are consistent with results in the L and S bands (Keller et al. 2025), giving us confidence in the sensitivity of our X-band search. With our procedure optimized for the X-band dataset, we are now conducting the complete analysis across the C and X bands. In addition to confirming that our analysis is robust against uncertainties in solar velocity and dark matter virial velocity, we hope to evaluate how different dark matter halo models affect our exclusion limits.
}

\newpage
